%%%%%%%%%%%%%%%%%%%- HXN
\def\sode{9}
\begin{name}
	{\tenchude}
	{\tendethi}
	{\tentruong}
	{\thoigian}
\end{name}

\caulc
\Opensolutionfile{ans}[ans/ans-HXN-\sode-T]
\begin{ex}%Câu 1
    Cho hàm số $y=f(x)$ có bảng biến thiên như sau:\\
    \centerline{\begin{tikzpicture}[>=stealth]
            \tkzTabInit[nocadre=false,lgt=1,espcl=2.5,deltacl=0.5]{$x$/.7 ,$y'$/.7,$y$/2}
            {$-\infty$ , $-2$ , $3$ , $+\infty$}
            \tkzTabLine{ , + , $0$ , - , $0$ , + , }
            \tkzTabVar{-/$-\infty$ , +/$4$ , -/$-3$ , +/$+\infty$}
    \end{tikzpicture}}
    Tổng giá trị cực đại và cực tiểu của hàm số $y=f(x)$ bằng
    \choice
    {\True $1$}
    {$-3$}
    {$4$}
    {$2$}
\end{ex}

\begin{ex}%Câu 2
    Trong không gian $Oxyz$ , cho mặt cầu $(S)$ có tâm $I(-1;2;1)$ và đi qua điểm $M(3;-1;4)$ . Phương trình của mặt cầu $(S)$ là:
    \choice
    {\True $(x+1)^2+(y-2)^2+(z-1)^2=34$}
    {$(x-1)^2+(y-2)^2+(z+1)^2=16$}
    {$(x+1)^2+(y-2)^2+(z-1)^2=16$}
    {$(x-1)^2+(y+2)^2+(z+1)^2=34$}
\end{ex}

\begin{ex}%Câu 3
    Nếu $\mathop{\int}_0^1f(x)\text{d}x=2$ và $\mathop{\int}_1^0g(x)\text{d}x=-3$ thì $\mathop{\int}_0^1[3f(x)-4g(x)]\text{d}x$ là:
    \choice
    {$6$}
    {\True $-6$}
    {$18$}
    {$12$}
\end{ex}

\begin{ex}%Câu 4
    Cho hình lăng trụ đứng $ABC.A' B' C'$ có đáy là tam giác vuông cân tại $B$ và $AB=6$ . Tính khoảng cách từ điểm $C$ đến mặt phẳng $(ABB'A')$
    \choice
    {$3$}
    {$6\sqrt{2}$}
    {$3\sqrt{2}$}
    {\True $6$}
\end{ex}

\begin{ex}%Câu 5
Đường tiệm cận xiên của đồ thị hàm số $ f(x)=\dfrac{2x^2-3x+1}{x+1}$ có phương trình là
\choice
{\True $y=2x-5$}
{$y=2x+5$}
{$y=2x-3$}
{$y=2x+3$}
\end{ex}

\begin{ex}%Câu 6
Trong không gian $ Oxyz$, cho vectơ $\vec{a}=(-1;0;2)$. Mệnh đề nào dưới đây đúng?
\choice
{$2\vec{a}=(2;0;-4)$}
{$2\vec{a}=(-2;0;-4)$}
{\True $2\vec{a}=(-2;0;4)$}
{$2\vec{a}=(2;0;4)$}
\end{ex}

\begin{ex}%Câu 7
Cho hàm số $ y=f(x)$ có đạo hàm $f'(x)=(x-2)(x+1)$, $\forall x\in\mathbb{R}$. Mệnh đề nào dưới đây đúng?
\choice
{Hàm số đồng biến trên $(-1;+\infty)$}
{Hàm số đồng biến trên $(-\infty;2)$}
{Hàm số nghịch biến trên $(-1;2)$}
{\True Hàm số nghịch biến trên $(-1;2)$}
\end{ex}

\begin{ex}%Câu 8
Trong không gian $ Oxyz,$ cho đường thẳng $ d\colon \dfrac{x-1}{2}=\dfrac{y+1}{1}=\dfrac{z}{-3}$. Mặt phẳng $(P)$ đi qua điểm $ A(1;0;1)$ và vuông góc với đường thẳng $d$ có phương trình là:
\choice
{\True $2x+y-3z+1=0$}
{$2x+y-3z-1=0$}
{$x+z+1=0$}
{$x+z-1=0$}
\end{ex}

\begin{ex}%Câu 9
Một người thống kê lại thời gian thực hiện các cuộc gọi điện thoại của người đó trong một tuần ở bảng sau: (đơn vị: giây)\\
\centerline{\begin{tblr}{|c|c|c|c|c|c|c|}
        \hline
        {Thời gian} & $[0;60)$ & $[60;120)$ & $[120;180)$ & $[180;240)$ & $[240;300)$ & $[300;360)$\\
        \hline
        Số cuộc gọi & $ 9$ & $ 9$ & $ 5$ & $ 7$ & $ 2$ & $ 1$\\
        \hline
\end{tblr}}\\
Khoảng tứ phân vị của mẫu số liệu ghép nhóm trên bằng
\choice
{$180$}
{\True $40$}
{$60$}
{$169$}
\end{ex}

\begin{ex}%Câu 10
Tập nghiệm $S$ của bất phương trình $\log_{\tfrac{1}{2}}\left(x-3\right)\ge\log_{\tfrac{1}{2}}4$ là
\choice
{\True $S=\left(3;7\right]$}
{$S=\left[3;7\right]$}
{$S=\left(-\infty;7\right]$}
{$S=\left[7;+\infty\right)$}
\end{ex}

\begin{ex}%Câu 11
Một công ty thống kê lương của nhân viên theo tuần (đơn vị: USD) theo bảng sau:\\
\centerline{
\begin{tblr}{
        colspec = {|c|c|c|c|c|c|},
    }
    \hline
    Lương theo tuần (USD) & [10; 20) & [20; 30) & [30; 40) & [40; 50) & [50; 60] \\\hline
    Số công nhân & 4 & 6 & 10 & 20 & 10 \\\hline
\end{tblr}
}
Độ lệch chuẩn của mẫu số liệu này bằng bao nhiêu (làm tròn tới hàng phần chục)?
\choice
{\True $11{,}7$}
{$12$}
{$11{,}4$}
{$12{,}5$}
\end{ex}

\begin{ex}%Câu 12
Cho hàm số $f(x)=x^2-\dfrac{4}{x}$. Giá trị của $\int\limits_1^2f'(x)\mathrm{d}x$ bằng
\choice
{$\dfrac{7}{3}-\text{ln}2$}
{$3$}
{$\dfrac{7}{3}$}
{\True $5$}
\end{ex}

\Closesolutionfile{ans}
\cauds
\Opensolutionfile{ans}[ans/ans-HXN-\sode-TF]

\begin{ex}%Câu 13
\immini
{
    Cho hàm số $y=f(x)$ xác định trên $\mathbb{R}\setminus\left[-1;1\right]$ liên tục trên mỗi khoảng xác định và có bảng biến thiên như sau
}
{
    \begin{tikzpicture}[>=stealth]
        \tkzTabInit[nocadre=false,lgt=1,espcl=2.5,deltacl=0.5]{$x$/.7 ,$y'$/.7,$y$/2}
        {$-\infty$ , $-1$,$1$ , $+\infty$}
        \tkzTabLine{ , +,t,h,t , - , }
        \tkzTabVar{-/$2$,+DH/$+\infty$, D+/$0$,-/$-2$}
    \end{tikzpicture}
}
\choiceTF
{Đồ thị hàm số $y=f(x)$ có đường tiệm cận đứng $x=1$}
{\True Đồ thị hàm số $y=f(x)$ có đúng hai đường tiệm cận ngang}
{\True Đồ thị hàm số $y=f(x)$ không có đường tiệm cận xiên}
{Đồ thị hàm số $y=\dfrac{1}{f(x)+1}$ có tất cả bốn đường tiệm cận}
\loigiai{
\begin{itemchoice}
    \itemch 
    \itemch 
    \itemch 
    \itemch Đặt $g(x)=\dfrac{1}{f(x)+1}$. \\
    Ta có $\lim\limits_{x\to -\infty }g(x)=\dfrac{1}{2+1}=\dfrac{1}{3}$ (vì $f(x)\to 2$); $\lim\limits_{x\to +\infty }g(x)=\dfrac{1}{-2+1}=-1$ (vì $f(x)\to -2$).\\
    Vì vậy đồ thị hàm số $y=g(x)$ có hai tiệm cận ngang $y=\dfrac{1}{3}$; $y=-1$.\\
    Xét $f(x)+1=0\Leftrightarrow f(x)=-1$. Phương trình này có một nghiệm thuộc khoảng $\left(1;+\infty \right)$.\\
    Do đó đồ thị hàm số $y=g(x)$ có một tiệm cận đứng.\\
    Vậy đồ thị hàm số $y=g(x)$ có tất cả ba đường tiệm cận.
\end{itemchoice}
}
\end{ex}

\begin{ex}%Câu 14
Có hai tên cướp vừa lấy được một chiếc ca nô ở vị trí $A$ thuộc bờ sông, chúng liền cho ca nô chạy theo phương hợp với bờ sông một góc $60^\circ$ với vận tốc $v=2t$ (mét/giây), trong đó $t$ (giây) là thời gian kể từ khi xuất phát. Sau $21$ giây, ca nô đến vị trí $B$ và chúng quyết định chuyển hướng cho ca nô chuyển động thẳng đều theo phương song song với bờ sông, tầm nửa phút sau thì ca nô đến $C$ (tham khảo hình vẽ).\\
\centerline{
    \includegraphics[width=8cm]{img/HXN-9-14}
}
    \choiceTF
    {\True Vị trí $B$ mà ca nô bọn cướp chuyển hướng cách bờ sông khoảng $382$ m (làm tròn đến hàng đơn vị)}
    {Khoảng cách $A$, $C$ tính theo đường chim bay bằng $1522$ m (làm tròn đến hàng đơn vị)}
    {\True Nếu các chiến sĩ công an khởi động ca nô và đi thẳng từ $D$ đến $C$ với vận tốc được tăng thêm $3$ m sau mỗi giây thì sau $21$ giây sẽ đến vị trí $D$ (làm tròn đến hàng đơn vị của giây)}
    {Trên thực tế các chiến sĩ đã chọn giải pháp cho ca nô khởi động và di chuyển vuông góc với bờ với gia tốc $a$ dương, cùng lúc đó bọn cướp từ vị trí $D$ tiến thẳng về phía trước (giữ nguyên hướng đi và tốc độ), hai bên giáp mặt nhau khi $a=4{,}78$ m/s$^2$ (làm tròn đến hàng phần trăm)}
    \loigiai{
        \begin{center}
            \begin{tikzpicture}[declare function={d=6;b=2;c=3;},thick]
                \path
                (0,0) coordinate (A)
                (-60:b) coordinate (B)
                ($(B)+(0:c)$) coordinate (C)
                (0:d) coordinate (D)
                ($(A)!(B)!(D)$) coordinate (H)
                ($(A)!(C)!(D)$) coordinate (K)
                ($(B)!(D)!(C)$) coordinate (E)
                pic[draw,angle radius=5mm,"$60^\circ$",angle eccentricity=1.7]{angle=B--A--H}
                ;
                \foreach \x/\y/\z in {B/H/K,C/K/D,C/E/D}\draw pic[draw,angle radius =2mm] {right angle = \x--\y--\z};
                \draw[dashed] (B)--(H) (C)--(K) (C)--(E)--(D);
                \draw (A)--(B)--(C) (D)--(A);
                \draw[red] (A)--(C)--(D);
                \foreach \x/\g in {A/180,B/-90,C/-90,D/0,E/-90,H/90,K/90}\draw[fill=white] (\x) circle (1pt)+(\g:3mm) node{$\x$};
            \end{tikzpicture}
        \end{center}
    \begin{itemchoice}
        \itemch Sau $21$ giây, ca nô bọn cướp đi được $AB=\int\limits_0^{21}{2t\mathrm{d}t}=441m$.\\
        Gọi $H$ là hình chiếu vuông góc của $B$ trên bờ sông, khi đó $BH=AB\sin 60^\circ=\dfrac{441\sqrt{3}}{2}\approx 382m$.
        \itemch Vận tốc của ca nô bọn cướp tại $B$ là $v_B=2\times 21=42$ m/s.\\
        Khoảng cách hai vị trí $B$, $C$ là $BC=42\times 30=1260m$.\\
        Khoảng cách hai vị trí $A$, $H$ là $AH=AB\cos 60^\circ=220{,}5m$.\\
        Gọi $K$ là hình chiếu vuông góc của $C$ trên bờ sông thì $HK=BC=1260m$.\\
        Do đó $AK=AH+HK=1480{,}5m$; $CK=BH=\dfrac{441\sqrt{3}}{2}m$ và $AC=\sqrt{AK^2+CK^2}\approx 1529 m$.
        \itemch Ta có $DK=2000-\left(220{,}5+1260\right)=519{,}5m$; suy ra $CD=\sqrt{CK^2+DK^2}\approx 644{,}8m$.\\
        Thời gian để các chiến sĩ đi được từ $D$ đến $C$ thỏa $\displaystyle\int\limits_0^t{3t\mathrm{d}t}=\sqrt{415741}\Rightarrow t\approx 21s$.
        \itemch Gọi $E$ là vị trí hai bên giáp mặt nhau (nếu có) thì tam giác $CDE$ vuông tại $E$.\\
        Khi đó $CE=DK=519{,}5m$ và $DE=\dfrac{441\sqrt{3}}{2}m$.\\
        Thời gian để ca nô bọn cướp đi từ $C$ đến $E$ là $\dfrac{CE}{42}=\dfrac{1039}{84}\approx 12{,}37s$.\\
        Gia tốc $a$ thỏa mãn $\displaystyle\int\limits_0^{\tfrac{1039}{84}}{at \mathrm{d}t}=\dfrac{441\sqrt{3}}{2}\Rightarrow 4{,}99m/s^2$.
    \end{itemchoice}
    }
\end{ex}

\begin{ex}%Câu 15
    \immini
    {
        Một chiếc máy bay thương mại Comac C919 đang bay trên bầu trời theo một đường thẳng từ $D$ đến $E$ có hình chiếu trên mặt đất là đoạn $CB$. Tại vị trí $D$ thì máy bay bay cách mặt đất $9000$m, tại vị trí $E$ thì máy bay cách mặt đất $12000$ m. Một ra đa được đặt trên mặt đất tại vị trí $O$ cách $C$ khoảng $20000$m, cách $B$ khoảng $16000$m và $\widehat{BOC}=90^\circ$; phạm vi theo dõi của ra đa là $20$ km.\\
        Xét hệ trục tọa độ $ Oxyz$ (đơn vị trên mỗi trục là $ 1000$ m) với $ O$ là vị trí đặt ra đa, $ B$ thuộc tia $Oy$, $C$ thuộc tia $ Ox$.
    }
    {
        \includegraphics[width=5cm]{img/HXN-9-15}
    }
    \choiceTF
    {Tại $ D$, máy bay cách ra đa $23000$m (làm tròn đến hàng nghìn theo đơn vị mét)}
    {\True Khi máy bay bay đến điểm $ I$ (m là trung điểm của $AB$), máy bay cách mặt đất $10500$ m}
    {\True Trên hành trình bay từ $ D$ đến $E$, máy bay sẽ đi qua điểm có tọa độ$ P(16;3,2;9,6)$}
    {Khoảng cách giữa vị trí đầu tiên và vị trí cuối cùng mà máy bay bay trong phạm vi theo dõi của ra đa là $22000$m (làm tròn đến hàng trăm theo đơn vị mét)}
    \loigiai{
        Tọa độ các điểm là$O(0;0;0)\,,B(0;16;0)\,,C(20;0;0)\,,D(20;0;9)\,,E(0;16;12)$.
        \begin{itemchoice}
            \itemch Ta có $\vec{OD}=(20;0;9) \Rightarrow OD=\sqrt{20^2+9^2}=\sqrt{481}$.\\
            Khi ở $D$, khoảng cách giữa máy bay và ra đa là $1000\times \sqrt{481}\approx 22\,000$ m.
            \itemch Tọa độ trung điểm của $DE$ là $I\left(10;8;\dfrac{21}{2}\right)$; cao độ điểm I là $\dfrac{21}{2}=10{,}5$.\\
            Vậy khi bay đến điểm $I$, máy bay cách mặt đất $10{,}5$ km hay $10\,500$m.
            \itemch Ta có $\vec{DE}=(-20;16;3)$. Phương trình đường thẳng $DE$ là $\heva{& x=20-20t \\& y=16t \\& z=9+3t } $ ($t$ là tham số thực).\\
            Thay tọa độ $P(16;3{,}2;9{,}6)$ vào phương trình $DE$ ta được $\heva{& 16=20-20t \\& 3{,}2=16t \\& 9{,}6=9+3t } \Leftrightarrow t=0{,}2\in (0;1)$.\\
            Do đó trên hành trình bay từ $D$ đến $E$, máy bay sẽ đi qua điểm $P(16;3{,}2;9{,}6)$.
            \itemch Gọi $H(20-20t;16t;9+3t)\in DE$ là hình chiếu vuông góc của $O$ trên $DE$.\\
            \centerline{\includegraphics[width=5cm]{img/HXN-9-15-LG}}
            Hai vectơ $\heva{& \vec{OH}=(20-20t;16t;9+3t) \\& \vec{DE}=(-20;16;3) } $ vuông góc với nhau nên\\
            $\vec{OH}\cdot \vec{DE}=0 \Leftrightarrow -20(20-20t)+16\cdot 16t+3(9+3t)=0\Leftrightarrow t=\dfrac{373}{665}$.\\
            Khi đó $\vec{OH}=\left(\dfrac{1\,168}{133};\dfrac{5\,968}{665};\dfrac{7\,104}{665}\right)$ và $OH=\dfrac{16\sqrt{469490}}{665}$.\\
            Gọi $M$, $N$ lần lượt là điểm đầu tiên và điểm cuối cùng mà máy bay xuất hiện trong phạm vi theo dõi của ra đa, ta có $MN=2\sqrt{R^2-OH^2}=2\sqrt{20^2-\dfrac{180736}{665}}=\dfrac{584\sqrt{665}}{665}\approx 22\,600$m.
        \end{itemchoice}
    }
\end{ex}


\begin{ex}%Câu 16
\immini
{
    Biết rằng $5\%$ cá heo trong một khu bảo tồn mắc một bệnh nhất định. Một xét nghiệm chẩn đoán có thể được áp dụng để xác định cá heo có mắc bệnh hay không:
\begin{itemize}
    \item Nếu cá heo mắc bệnh, xét nghiệm cho kết quả dương tính với xác suất $0{,}96$.
    \item Nếu cá heo không mắc bệnh, xét nghiệm cho kết quả dương tính với xác suất $0{,}02$.
\end{itemize}
Xét nghiệm được áp dụng cho một pé cá heo được chọn ngẫu nhiên.
}
{
    \includegraphics[width=5cm]{img/HXN-9-16}
}
    \choiceTF
    {Xác suất để thu được kết quả dương tính bằng $0{,}07$}
    {Biết rằng kết quả dương tính thu được, xác suất để cá heo này thực sự mắc bệnh bằng $\dfrac{45}{67}$}
    {\True Cá heo khi xét nghiệm đã cho kết quả dương tính lần đầu, xác suất để khi xét nghiệm lần tiếp theo vẫn cho kết quả dương tính bằng $ 0{,}69$ (làm tròn đến hàng phần trăm)}
    {Biết rằng kết quả xét nghiệm lần thứ hai là dương tính, xác suất để cá heo này thực sự mắc bệnh bằng $0{,}97$ (làm tròn đến hàng phần trăm)}
    \loigiai{
        Gọi $A$ là biến cố: \lq\lq Cá heo thực sự mắc một bệnh nhất định\rq\rq.\\
        $B_i$ là biến cố: \lq\lq Xét nghiệm cá heo cho kết quả dương tính lần thứ $i$\rq\rq; với $i\in \mathbb{N}^*$.
    \begin{itemchoice}
        \itemch Xác suất để có kết quả dương tính là $P\left(B_1\right)=P(A)\cdot P\left(B_1\mid A\right)+P\left({\bar{A}}\right)\cdot P\left(B_1\mid \bar{A}\right)$;\\
        $P\left(B_1\right)=0{,}05\times 0{,}96+0{,}95\times 0{,}02=0{,}067$.
        \itemch Ta có $P\left(A|B_1\right)=\dfrac{P(AB)}{P(B)}=\dfrac{0{,}05\times 0{,}96}{0{,}067}=\dfrac{48}{67}$.
        \itemch Ta có $P\left(B_2|B_1\right)=\dfrac{P\left(B_1B_2\right)}{P\left(B_1\right)}=\dfrac{0{,}05\times 0{,}96^2+0{,}95\times 0{,}02^2}{0{,}067}\approx 0{,}69$.
        \itemch Ta có: $P\left(A|B_2\right)=\dfrac{P\left(AB_2\right)}{P\left(B_2\right)}=\dfrac{0{,}05\times 0{,}96^2}{0{,}05\times 0{,}96^2+0{,}95\times 0{,}02^2}=\dfrac{2304}{2323}\approx 0{,}99$.
    \end{itemchoice}
    }
\end{ex}

\Closesolutionfile{ans}
\caukq
\Opensolutionfile{ans}[ans/ans-HXN-\sode-SA]

\begin{ex}%Câu 17
    \immini
    {
        Trong một trò chơi, người chơi muốn tìm đường đi ngắn nhất để đi từ A đến P, biết rằng từ A đến P có những đường đi được cho như hình vẽ; thông số trên các đoạn thẳng chính là khoảng cách giữa hai vị trí tương ứng. Đường đi thoả mãn điều kiện trên nhận giá trị nhỏ nhất là bao nhiêu?
    \shortans{21}
    }
    {
        \includegraphics[width=7cm]{img/HXN-9-17}
    }
    \end{ex}
    
    \begin{ex}%Câu 18
\immini
{
    Một vật chuyển động theo quy luật $ s=s(t)=\dfrac{1}{3}{t^3}-\dfrac{3}{2}{t^2}+10t+2$; với $ t$(giây) là khoảng thời gian tính từ lúc vật bắt đầu chuyển động và $ s$(mét) là quãng đường vật đi được trong thời gian đó). Tính quãng đường mà vật đi được khi vận tốc nó đạt $ 20$m/s (kết quả làm tròn đến hàng phần chục).
\shortans{54,2}
}
{
    \includegraphics[width=5cm]{img/HXN-9-18}
}
\end{ex}

\begin{ex}%Câu 19
\immini
{
    Một cái hộp đựng đồ chơi có dạng hình lập phương $ ABCD.A'B'C'D'$. Gọi $\varphi$ là góc giữa hai mặt phẳng $\left(AB'C'\right)$ và $\left(AC'D'\right)$, tính giá trị $\cot\varphi$ và làm tròn đến hàng phần trăm.
\shortans{0,58}
}
{
    \includegraphics[width=5cm]{img/HXN-9-19}
}
\loigiai{
\immini
{
    \textbf{Nhận xét:} $\left(AB'C'\right)\equiv \left(ADC'B'\right)$, $\left(AC'D'\right)\equiv \left(ABC'D'\right)$.\\
    Ta có: $\heva{& CD'\perp C'D \\& CD'\perp AD\left(doAD\perp \left(CDD'C'\right)\right) } $\\
    $\Rightarrow CD'\perp \left(ADC'B'\right)$\tagEX{1}
    Tương tự: $\heva{& B'C\perp BC' \\& B'C\perp AB\left(doAB\perp \left(BCC'B'\right)\right) } $\\
    $\Rightarrow B'C\perp \left(ABC'D'\right)$\tagEX{2}
    Từ $(1)$ và $(2)$ suy ra $\left(\left(AB'C'\right),\left(AC'D'\right)\right)=\left(CD',CB'\right)$.\\
    Giả sử cạnh hình lập phương bằng $a$.\\
    Ta có $CB'=CD'=B'D'=a\sqrt{2}$ (đường chéo trong hình vuông). Suy ra tam giác $CB'D'$ đều.\\
    Do vậy $\varphi =\left(\left(AB'C'\right),\left(AC'D'\right)\right)=\left(CD',CB'\right)=\widehat{B'CD'}=60^\circ$\\
    Vậy $\cot \varphi \approx 0{,}58$.
}
{
    \includegraphics[width=5cm]{img/HXN-9-19-LG}
}
}
\end{ex}

\begin{ex}%Câu 20
\immini
{
    Trong một lễ hội mùa hè, có một trò chơi mà mỗi lần chơi, người chơi sẽ tung đồng thời bốn đồng xu cân đối một cách ngẫu nhiên. Người chơi chỉ thắng cuộc nếu nhận được ít nhất ba mặt ngửa từ bốn đồng xu đã tung.
Xác suất để trong $5$ lần chơi, người chơi thắng được ít nhất ba lần xấp xỉ bằng $ m\cdot 10^{-2}$ với $ m$ là số tự nhiên có hai chữ số (đã được làm tròn đến hàng đơn vị); hỏi giá trị của $m$ bằng bao nhiêu?
\shortans{18}
}
{
\includegraphics[width=5cm]{img/HXN-9-20}
}
\loigiai{
\textbf{Bước 1:} Tính xác suất thắng trong một lần chơi\\
Trong một lần chơi, người chơi tung $4$ đồng xu cân đối:
\begin{itemize}
    \item Xác suất nhận được đúng 3 mặt ngửa là $\mathrm{C}_4^3\cdot \left(\dfrac{1}{2}\right)^3\cdot \left(\dfrac{1}{2}\right)^1=4\cdot \dfrac{1}{8}\cdot \dfrac{1}{2}=4\cdot \dfrac{1}{16}=\dfrac{1}{4}$.
    \item Xác suất nhận được cả 4 mặt ngửa là $\left(\dfrac{1}{2}\right)^4=\dfrac{1}{16}$.\\
    Vậy xác suất thắng trong một lần chơi là: $\dfrac{1}{4}+\dfrac{1}{16}=\dfrac{5}{16}$;\\
    xác suất thua trong một lần chơi là $1-\dfrac{5}{16}=\dfrac{11}{16}$.
\end{itemize}
\textbf{Bước 2:} Tính xác suất thắng ít nhất $3$ lần trong $5$ lần chơi.\\
Xác suất này bao gồm các trường hợp:
\begin{itemize}
    \item Thắng đúng 3 lần, thua 2 lần: $\mathrm{C}_5^3\cdot \left(\dfrac{5}{16}\right)^3\cdot \left(\dfrac{11}{16}\right)^2=10\cdot \dfrac{125}{4096}\cdot \dfrac{121}{256}=\dfrac{151250}{1048576}$.
    \item Thắng đúng 4 lần, thua 1 lần: $\mathrm{C}_5^4\cdot \left(\dfrac{5}{16}\right)^4\cdot \left(\dfrac{11}{16}\right)^1=5\cdot \dfrac{625}{65536}\cdot \dfrac{11}{16}=\dfrac{34375}{1048576}$.
    \item Thắng cả 5 lần: $\left(\dfrac{5}{16}\right)^5=\dfrac{3125}{1048576}$.\\
    Xác suất cần tính là: $\dfrac{151250}{1048576}+\dfrac{34375}{1048576}+\dfrac{3125}{1048576}\approx 18\cdot 10^{-2}$.
\end{itemize}
}
\end{ex}

\begin{ex}%Câu 21
\immini
{
    Một nghệ nhân muốn thiết kế một chiếc ly hình nón có thể chứa được tối đa $500$ ml nước với bán kính miệng ly $ R<8cm$. Sau đó, anh ta đặt một khối trụ đặc có bán kính đáy $2$ cm vào bên trong ly sao cho trục của hình trụ trùng với trục của hình nón và đáy trên của hình trụ nằm cùng một mặt phẳng với đáy của hình nón. Nghệ nhân muốn đặt các thanh thủy tinh phát sáng có dạng đoạn thẳng vào ly chứa đầy dung dịch màu. Biết rằng $AB$ là thanh thủy tinh có độ dài lớn nhất có thể đặt vào ly (không có điểm nào nhô ra khỏi mặt nước), tìm giá trị lớn nhất đó theo đơn vị cm (làm tròn đến hàng phần mười).
\shortans{9,5}
}
{
\begin{tikzpicture}[scale=1, font=\footnotesize, thick,declare function={
        H=-7;R=5;
        k=.8/2;
        h=k*H;r=k*R;
        c=1/6;
        G=asin(R*c/H);
        g=asin(r*c/h);
        X=R*cos(G);Y=R*sin(G);
        xd=r*cos(g);yd=r*sin(g);
    },line join=round, line cap=round,x=.8cm]
    \path
    (X,Y) coordinate (M)
    (xd,yd) coordinate (m)
    ($(m)+(90:{H-h})$) coordinate (m1)
    (90:H)coordinate (S)
    (m) arc (g:180-g:{r} and {r*c}) coordinate (ma)
    ($(ma)+(90:{H-h})$) coordinate (m1a)
    ;
    \draw[dashed]
    (ma)--(m1a)
    (m)--(m1) arc (g:180-g:{r} and {r*c})
    (m1) arc (360+g:180-g:{r} and {r*c})
    ;
    \draw 
    (m) arc (g:360-g:{r} and {r*c})
    (M) arc (G:180-G:{R} and {R*c})
    (S)--(M) arc(360+G:180-G:{R} and {R*c})--cycle;
    \draw[<->] ([shift={(0:R+.2)}]S)--++(-90:H)node[midway,fill=white]{$5$ cm};
\end{tikzpicture}
}
\loigiai{
    \begin{center}
        \begin{tikzpicture}[scale=1, font=\footnotesize, thick,>=stealth, line join=round, line cap=round, declare function={
                H=-7;R=5;
                k=2/5;
                h=k*H;r=k*R;
                c=1/5;
                G=asin(R*c/H);
                g=asin(r*c/h);
                X=R*cos(G);Y=R*sin(G);
                xd=r*cos(g);yd=r*sin(g);
            },x=.8cm]
            \path
            (0,0) coordinate (O)
            (X,Y) coordinate (M)
            (xd,yd) coordinate (m)
            ($(m)+(90:{H-h})$) coordinate (m1)
            (90:H)coordinate (S)
            (m) arc (g:180-g:{r} and {r*c}) coordinate (ma)
            ($(ma)+(90:{H-h})$) coordinate (m1a)
            ;
            \draw[dashed]
            (ma)--(m1a)
            (m)--(m1) arc (g:360-g:{r} and {r*c})
            ;
            \path
            (M)arc (G:20:{R} and {R*c}) coordinate (B)
            (m) arc (g:-70:{r} and {r*c}) coordinate (C)
            (m1)arc (g:-70:{r} and {r*c}) coordinate (A)
            pic[draw,angle radius=3mm]{right angle=O--C--B};
            \draw 
            (m) arc (g:360-g:{r} and {r*c})
            (M) arc (G:360-G:{R} and {R*c})
            (S)--(M) arc(360+G:180-G:{R} and {R*c})--cycle;
            \draw[<->] ([shift={(180:R+.2)}]S)--++(-90:H)node[midway,fill=white]{$5$ cm};
            \draw[dashed] (C)--(A)--(B) (O)--(90:{H-h})node[midway,fill=white]{$h_t$};
            \draw[red] (O)--(C)node[midway,left,black]{$r$}--(B)--(O)node[midway,above,black]{$R$};
            \foreach \x/\g in {A/-90,B/30,C/-135,O/-170}\draw[fill=white] (\x) circle (1pt)+(\g:3mm) node{$\x$};
        \end{tikzpicture}
    \end{center}
Thể tích khối nón: $V=\dfrac{1}{3}\pi R^2h=500ml=500cm^3 \Rightarrow h=\dfrac{1500}{\pi R^2}$. \tagEX{1}
Gọi $h_t$ là chiều cao hình trụ, $r=2cm$ là bán kính đáy hình trụ.\\
Do hình trụ nằm trong hình nón và có cùng trục, ta có tỉ lệ: $\dfrac{r}{R}=\dfrac{h-h_t}{h}\Rightarrow h_t=h\left(1-\dfrac{r}{R}\right)$. \tagEX{2}
Thay $(1)$ vào $(2)$: $h_t=\dfrac{1500}{\pi R^2}\left(1-\dfrac{r}{R}\right)$. \tagEX{3}
Thanh $AB$ nằm chéo trong không gian giữa hình nón và hình trụ.\\
Gọi $A$ là điểm nằm dưới (thuộc đường tròn đáy dưới hình trụ), $B$ nằm trên mặt nước. Khi đó, độ dài $AB$ được tính bằng công thức: $AB=\sqrt{R^2-4+h_t^2}$. \tagEX{4}
Thay $(3)$ vào $(4)$: $AB=\sqrt{R^2-4+\left(\dfrac{1500}{\pi R^2}\left(1-\dfrac{2}{R}\right)\right)^2}=f(R)$.\\
Khảo sát hàm $f(R)$, ta có $f(R)=AB$ đạt giá trị lớn nhất xấp xỉ $9{,}5cm$; khi đó $R\approx 7{,}7$.
}
\end{ex}

\begin{ex}%Câu 22
Trong không gian $ Oxyz$, cho hai điểm $ A(10;6;-2)$, $B(5;10;-9)$ và mặt phẳng $(\alpha)\colon 2x+2y+z-12=0$. Điểm $ M$ di động trên $\left(\alpha\right)$ sao cho $ MA$, $ MB$ luôn tạo với $(\alpha)$ các góc bằng nhau. Biết rằng $ M$ luôn thuộc một đường tròn $(C)$ cố định có tâm $ H(a;b;c)$. Tính tổng $a^2+b^2+c^2$.
\shortans{248}
\loigiai
{
\immini
{
    Gọi $H$, $K$ lần lượt là hình chiếu vuông góc của $A$, $B$ trên mặt phẳng $\left(\alpha \right)$.\\
Khi đó $AH=d\left(A,\left(\alpha \right)\right)=6$; $BK=d\left(B,(\alpha )\right)=3$.\\
Vì $MA$, $MB$ tạo với $\left(\alpha \right)$ các góc bằng nhau nên $\widehat{AMH}=\widehat{BMK}$ mà $AH=2BK$ suy ra $MA=2MB$.\\
Gọi $M(x;y;z)$, ta có $MA=2MB\Leftrightarrow MA^2=4MB^2$\\
$ \Leftrightarrow x^2+y^2+z^2-\dfrac{20}{3}x-\dfrac{68}{3}y+\dfrac{68}{3}z+228=0$.
}
{
    \includegraphics[width=5cm]{img/HXN-9-22-LG-a}
}
\immini
{
    Như vậy, điểm $M$ nằm trên mặt cầu $(S)$ có tâm $I\left(\dfrac{10}{3};\dfrac{34}{3};-\dfrac{34}{3}\right)$ và bán kính $R=2\sqrt{10}$.\\
Mặt khác ta có $M$ di động trên $\left(\alpha \right)$, vì vậy tập hợp điểm $M$ chính là đường tròn giao tuyến $(C)$ được tạo bởi mặt cầu $(S)$ và mặt phẳng $\left(\alpha \right)$.\\
Gọi $H$ là tâm của đường tròn $(C)$, khi đó $H$ là hình chiếu vuông góc của $I$ trên mặt phẳng $\left(\alpha \right)$.\\
Đường thẳng $d$ qua $I$ và vuông góc  $\left(\alpha \right)$ có phương trình  $d\colon \heva{& x=\dfrac{10}{3}+2t \\& y=\dfrac{34}{3}+2t \\& z=-\dfrac{34}{3}+t }$.
}
{
    \includegraphics[width=5cm]{img/HXN-9-22-LG-b}
}
Thay phương trình  $d$ vào $\left(\alpha \right)\colon 2\left(\dfrac{10}{3}+2t\right)+2\left(\dfrac{34}{3}+2t\right)+\left(-\dfrac{34}{3}+t\right)-12=0\Leftrightarrow t=-\dfrac{2}{3}$.\\
Từ đó suy ra $H(2;10;-12)$ và $a=2$, $b=10$, $c=-12\Rightarrow a^2+b^2+c^2=248$.
}
\end{ex}

\Closesolutionfile{ans}
\inputansbox{6,4,3}{ans/ans-HXN-\sode-T,ans/ans-HXN-\sode-TF,ans/ans-HXN-\sode-SA}
